\chapter*{Abstract}

RISC-V is an open instruction set architecture (ISA) introduced in 2010 at the University of California, Berkeley as a royalty-free alternative to the proprietary instruction sets that dominated the industry. It was originally conceived as an academic project to support research in computer architecture. Nevertheless, its open and standardised nature has contributed to its widespread adoption and has established it as a viable alternative to traditional architectures such as x86 and ARM in a variety of application domains.

This document summarises the process of implementating of a five-stage pipelined RISC-V processor core, capable of executing the base instruction set defined by the standard.

The project is based on a single-cycle design developed at National Cheng Kung University (Taiwan). This design has been thoroughly analysed and serves as the foundation for the work presented herein.

Finally, the document describes the deployment of the resulting design on an FPGA development board using the Vivado design suite.\newline\newline

\textbf{Palabras clave:} RISC-V, Chisel, HDL, computer architecture, hardware design.